\section{Experimental Setup}
\label{sec:setup}

% Draft for Ammaar's revision. Week-2 deliverable. Values here mirror
% configs/default.yaml and the training/eval scripts — the numbers
% audit (Phase G) must cross-check this section against the config.

\subsection{Splits, Seeds, and Reproducibility}
All experiments use the frozen 895/192/\NumTestImages{} split
(Sec.~\ref{sec:dataset}); every reported number is on the
\NumTestImages-image test split against human-verified GT only. Every
run writes a metadata record (full configuration, git SHA, seed,
library versions), and a global seed fixture makes pipeline runs
deterministic. The detector is trained with three seeds; pipeline
results are reported for each detector seed (mean$\pm$std), and
bootstrap confidence intervals capture per-image variation.

\subsection{Training and Hardware}
The YOLOv8s detector \cite{yolov8} is trained at $640$\,px on the
train split with the val split for model selection (an NVIDIA RTX~3090
was used for training; \todoa{confirm final training hardware +
epochs/hyperparameters from experiments/train\_all/ before
submission}). Everything downstream — the full pipeline, benchmark,
and ngspice simulation — runs locally on an Apple~M1 laptop (CPU +
Metal inference), which is part of the design brief: a lightweight
pipeline that does not require server-side inference.

\subsection{Metric Cascade (C4)}
\label{sec:metrics}
We report, in order of increasing strictness:
\begin{itemize}
  \item \textbf{Detection}: mAP@0.5 and mAP@0.5:0.95 with per-class AP
  and supports.
  \item \textbf{Terminal-pair $F_1$}: over all unordered pairs of
  component terminals, whether the pair shares a net in prediction vs.\
  GT — the finest-grained connectivity signal.
  \item \textbf{Net-level $F_1$}: Hungarian matching between predicted
  and GT nets (a net = a set of terminals), scored by set overlap.
  \item \textbf{Per-component connected accuracy}: fraction of
  components with \emph{every} terminal on the correct net.
  \item \textbf{Normalized graph edit distance} (nGED): between
  predicted and GT terminal graphs, with a timeout and documented
  upper-bound fallback on large graphs.
  \item \textbf{Strict end-to-end success}: an image counts only if
  every component is detected, classified, and fully correctly
  connected — a product over $\sim$14 components per image, so it is
  deliberately punishing.
  \item \textbf{SPICE validity and DC solvability}: whether ngspice
  \cite{ngspice} parses the netlist, and whether DC operating-point
  analysis converges — reported \emph{before and after} repair (C5),
  on a separate axis from topology.
\end{itemize}

\textbf{Fair alignment.} Predicted components are aligned to GT
components by Hungarian matching on bounding-box IoU within class
(threshold 0.3); unmatched components on either side are retained with
disjoint identities so they \emph{penalize} the graph metrics rather
than being dropped. Terminal order of symmetric parts is canonicalized
by an identical connectivity signature on both sides, so no metric
rewards or punishes arbitrary bookkeeping.

Alignment depends on ground-truth box geometry, which — unlike the net
assignments — was not human-verified (Sec.~\ref{sec:gt}). We therefore
rebuild GT box geometry from the published COCO annotations before
scoring and report the sensitivity of the results to the alignment
threshold, rather than choosing a threshold that flatters the numbers.
\draftnote{finalize once the corrected-box comparison run lands; state
which GT the headline table uses and show the sensitivity curve.}

\textbf{Statistics.} Aggregate metrics carry bootstrap 95\% CIs
(per-image resampling). Configuration comparisons on the same images
are \emph{paired}: significance comes from the bootstrap distribution
of per-image deltas (2{,}000 resamples), never from overlap of two
independent CIs.

\subsection{Oracle Stage Attribution}
The GT-injection oracle replaces one pipeline stage at a time with its
ground-truth output (perfect detections; GT-rendered wires; oracle
snapping) and measures the end-to-end gain, attributing the error
budget across stages. \draftnote{describe mode C rendering here once
the orthogonal re-render is validated.}

\subsection{Port Templates (C3)}
Port templates are fitted on the Digitize-HCD port-location corpus with
a deterministic held-out split (every fifth crop), and evaluated by
normalized port-localization error on the held-out crops. All
\NumClasses{} classes have templates; note that several classes' ports
are unnamed upstream (both AC sources and the one-port DC source), so
for those the template fixes an ordering convention only and no
polarity claim is made.

\subsection{Repair Evaluation (C5)}
Repair is evaluated on its own axis: DC-solvability rate before vs.\
after repair, assumptions per circuit (gauge vs.\ assumption split),
and the topology-preservation check (topology metrics bit-identical
with repair on/off). \todoa{[IDEAL] expert-acceptance study: mentor
rates $\sim$30 ledgers correct-minimal / over-repaired / wrong.}
